\section{Review Methodology}
\label{sec:method}

This study uses a structured, chronological narrative review of legal artificial intelligence.
Instead of grouping the literature by topic, it follows the order of publication. Each advance
appears after the work that motivated it and before the work that responded to it. This design
suits a field that spans several related but distinct areas. These include legal knowledge
representation, natural language processing, information retrieval, large language models, and
evidence verification. The design also matches the review's purpose directly. We want to explain how
each change in system design responded to a limitation of what came before it, not to rank models by
a single reported score.

We identified relevant literature across computer science and legal informatics sources, including
Scopus, Web of Science, IEEE Xplore, the ACM Digital Library, the ACL Anthology, Google Scholar,
official publisher and proceedings pages, and arXiv. The search combined four term groups: the legal
domain (for example, ``legal AI'' or ``case-based reasoning''), the underlying technology (``legal
language model'', ``LEGAL-BERT''), retrieval and document structure (``hybrid retrieval'',
``structure-aware indexing'', ``knowledge graph''), and reliability (``legal hallucination'',
``citation verification''). We supplemented keyword search with backward and forward citation
searching to catch studies that different terminology might have hidden. Table~\ref{tab:sel}
summarizes the selection criteria we applied to the resulting records.

\begin{table*}[t]
\centering
\caption{Study selection criteria.}
\label{tab:sel}
\footnotesize
\setlength{\tabcolsep}{5pt}
\renewcommand{\arraystretch}{1.25}
\begin{tabularx}{\textwidth}{>{\raggedright\arraybackslash}p{2.6cm} >{\raggedright\arraybackslash}X >{\raggedright\arraybackslash}X}
\toprule
\textbf{Criterion} & \textbf{Include} & \textbf{Exclude} \\
\midrule
Topical relevance & Legal reasoning, legal NLP, retrieval, generation, verification, or evaluation. & AI in an unrelated domain with no meaningful link to legal information processing. \\
Technical contribution & A method, architecture, dataset, benchmark, or substantive empirical finding. & Mentions legal AI only in passing, with no relevant technical content. \\
Venue quality & A peer-reviewed article or proceedings paper in an established, well-indexed venue (a Q1-ranked journal or a recognized conference); foundational works; preprints that add substantial, clearly identified value. & Venues with no established standing in computer science or legal informatics, and non-scholarly material used as if it were primary research. \\
Evidence availability & An accessible bibliographic record, checked against the publisher's page or official proceedings, with a verifiable DOI where one exists. & Claims or publication details that cannot be independently verified. \\
Jurisdiction and language & Any jurisdiction or language relevant to the review's research questions. & No jurisdiction or language excluded solely for being under-represented. \\
\bottomrule
\end{tabularx}
\end{table*}

We included a record only after checking its bibliographic details against the publisher's page,
the official proceedings, or another authoritative source, and after verifying its DOI where one
exists. For preprints, we treated the consulted version as distinct from any later published
version. We checked venue quality deliberately. This survey keeps to Q1-ranked journals and
recognized conference venues, and we dropped a number of candidate papers at this stage because we
could not confirm their venue met that bar. For each included study, we recorded its research
objective, legal domain and jurisdiction, technical approach, evidence sources, evaluation method,
and acknowledged limitations. We then sorted this information into the thematic categories the
following sections use: historical foundations, legal language models, retrieval architecture,
document and knowledge representation, generation and orchestration, reliability and verification,
and evaluation design.

We present the resulting set of studies in original publication order, year by year, instead of
grouping them by theme. This lets the historical account in Section~\ref{sec:foundations} and onward read as one
continuous development rather than a set of separately ordered topics. We report quantitative
results only together with their original experimental context, because studies differ in
jurisdiction, corpus, query set, and definition of success. A legal classification benchmark, for
instance, cannot be compared directly with a human evaluation of unsupported legal claims. Throughout
the survey we keep three kinds of statements separate: findings individual studies report directly,
interpretations we draw by comparing multiple studies, and open questions for future research. That
discipline, together with our restriction to independently verifiable, well-indexed sources, keeps
this a critical narrative of the field rather than a catalogue of publications.
